% Generated from tex/chapters/ch04/08_構築技術.tex for the SWEBOK structure tree
\subsubsection{実行時設定と国際化}

実行時設定とは、プログラムの値や動作を、実行時に設定ファイルや環境変数などから読み取って決めることである。これにより、コードを変更せずに動作を変えられる。設定を使うことで柔軟性は高まるが、設定値の検証、既定値、変更履歴、機密情報の扱いが重要になる。

国際化とは、複数の言語や地域に対応できるようにソフトウェアを準備することである。これに対応して、特定の地域や言語に合わせる作業を地域化という。画面の文言、エラーメッセージ、日付形式、通貨、文字コード、並び順、右から左へ書く言語などを考慮する必要がある。

国際化を後から追加すると、コードのあちこちに埋め込まれた文字列や地域依存の処理を修正する必要がある。構築時から、文字列を外部化し、文字コードとロケールを意識することが重要である。
